\section*{Capítulo \ref{ch:b1:rock-forming-minerals} --- Minerais formadores de rochas}

\begin{solution}{exo:b1:rock-forming-minerals:1}
\ce{SiO4^{4-}}: \omterm{def:b1:rock-forming-minerals:silicate-classes}{nesossilicato}, $s = 0$. \ce{Si2O7^{6-}}: \omterm{def:b1:rock-forming-minerals:silicate-classes}{sorossilicato},
$s = 1$. \ce{Si6O18^{12-}}: \omterm{def:b1:rock-forming-minerals:silicate-classes}{ciclossilicato}, $s = 2$. \ce{Si2O5^{2-}}:
\omterm{def:b1:rock-forming-minerals:silicate-classes}{filossilicato}, $s = 3$. \ce{SiO2}: \omterm{def:b1:rock-forming-minerals:silicate-classes}{tectossilicato}, $s = 4$. \ce{SiO3^{2-}}
tem $4 - s/2 = 3$, logo $s = 2$: um anel (\omterm{def:b1:rock-forming-minerals:silicate-classes}{ciclossilicato}) ou uma cadeia simples
(\omterm{def:b1:rock-forming-minerals:silicate-classes}{inossilicato}); a razão sozinha não os distingue.
\end{solution}

\begin{solution}{exo:b1:rock-forming-minerals:2}
Forsterita: $2\times24.305 + 28.085 + 4\times15.999 =
\qty{140.691}{g.mol^{-1}}$; fayalita: $2\times55.845 + 28.085 +
4\times15.999 = \qty{203.771}{g.mol^{-1}}$. Volumes molares:
$140.691/3.275 = \qty{42.96}{cm^3.mol^{-1}}$ e $203.771/4.392 =
\qty{46.40}{cm^3.mol^{-1}}$.
\end{solution}

\begin{solution}{exo:b1:rock-forming-minerals:3}
$(0.78 - 0.72)/0.72 = +\qty{8.3}{\percent}$ para \ce{Fe^{2+}};
$-\qty{4.2}{\percent}$ para \ce{Ni^{2+}}; $+\qty{15.3}{\percent}$ para
\ce{Mn^{2+}}; $+\qty{38.9}{\percent}$ para \ce{Ca^{2+}}. Ferro e níquel,
abaixo de \qty{15}{\percent}, substituem livremente o magnésio; o manganês, logo acima, é limitado; o cálcio, além de \qty{30}{\percent}, quase não entra. Os quatro têm a mesma carga do magnésio.
\end{solution}

\begin{solution}{exo:b1:rock-forming-minerals:4}
$(3.51 - 2.16)/2.16 = \qty{62.5}{\percent}$ do grafite para o diamante;
$(2.93 - 2.71)/2.71 = \qty{8.1}{\percent}$ da calcita para a aragonita. Os \omterm{def:b1:rock-forming-minerals:polymorph}{polimorfos} mais densos, o diamante e a aragonita, são as formas de alta pressão.
\end{solution}

\begin{solution}{exo:b1:rock-forming-minerals:5}
Com $x = 0.30$ em $\mathrm{Na}_{1-x}\mathrm{Ca}_x\mathrm{Al}_{1+x}\mathrm{Si}_{3-x}\mathrm{O}_8$:
$\mathrm{(Na_{0.7}Ca_{0.3})(Al_{1.3}Si_{2.7})O_8}$. Cargas:
$0.7 + 2\times0.3 + 3\times1.3 + 4\times2.7 = 0.7 + 0.6 + 3.9 + 10.8 = +16$,
e $8\times(-2) = -16$.
\end{solution}

\begin{solution}{exo:b1:rock-forming-minerals:6}
Um grupo que se repete tem dois tetraedros compartilhando três vértices (aqueles cujo ápice une as duas cadeias) e dois compartilhando dois. Os quatro tetraedros têm
$4\times4 = 16$ vértices, dos quais $3 + 3 + 2 + 2 = 10$ são compartilhados, isto é,
$10/2 = 5$ oxigênios compartilhados; os $16 - 10 = 6$ outros não são compartilhados. Total
$5 + 6 = 11$ oxigênios para 4 silícios: \ce{Si4O11}, com carga
$4\times4 - 11\times2 = -6$. A proposição dá o mesmo com
$s = 5/2$: $4 - 5/4 = 11/4$ oxigênios por silício, carga $5/2 - 4 = -3/2$ por
silício.
\end{solution}

\begin{solution}{exo:b1:rock-forming-minerals:7}
Moles em \qty{100}{g}: \ce{SiO2} $40.20/60.083 = 0.66907$; \ce{MgO}
$47.10/40.304 = 1.16862$; \ce{FeO} $12.20/71.844 = 0.16981$; \ce{NiO}
$0.40/74.692 = 0.00536$. Oxigênios: $2\times0.66907 + 1.16862 + 0.16981 +
0.00536 = 2.68194$; $f = 4/2.68194 = 1.49146$. Por quatro oxigênios: Si
\num{0.998}, Mg \num{1.743}, Fe \num{0.253}, Ni \num{0.008}. Os
sítios octaédricos têm $1.743 + 0.253 + 0.008 = 2.004$, o sítio tetraédrico
\num{0.998}: uma boa análise, \ce{(Mg_{1.743}Fe_{0.253}Ni_{0.008})SiO4}.
$\mathrm{Fo} = 100\times1.743/(1.743 + 0.253) = 87.3$. O fator $f$
multiplica Mg e Fe do mesmo modo, de modo que se cancela na razão:
$\mathrm{Fo} = 100\,(47.10/40.304)/(47.10/40.304 + 12.20/71.844) =
100\times1.16862/1.33843 = 87.3$.
\end{solution}

\begin{solution}{exo:b1:rock-forming-minerals:8}
Positivas: $1 + 2\times3 + 3\times4 + 3 = +22$; negativas:
$10\times(-2) + 2\times(-1) = -22$. Uma folha de composição \ce{Si4O10} teria carga $-4$ (\cref{prop:b1:rock-forming-minerals:si-o}, $s = 3$); trocar um silício em cada quatro por alumínio acrescenta uma carga negativa, logo \ce{(Si3Al)O10}
tem carga $-5$. Os dois alumínios octaédricos com as duas hidroxilas,
$6 - 2 = +4$, equilibram quatro das cinco; o íon potássio entre as folhas equilibra a quinta.
\end{solution}

\begin{solution}{exo:b1:rock-forming-minerals:9}
(a) $M = (140.691 + 203.771)/2 = \qty{172.231}{g.mol^{-1}}$, $V = (42.96 +
46.40)/2 = \qty{44.68}{cm^3.mol^{-1}}$, $\rho = 172.231/44.68 =
\qty{3.855}{g.cm^{-3}}$; a média das densidades é $(3.275 + 4.392)/2 =
\qty{3.834}{g.cm^{-3}}$, \qty{0.021}{g.cm^{-3}} a menos. (b) Resolva
$203.771 + x(140.691 - 203.771) = 3.600\,\bigl(46.40 + x(42.96 - 46.40)\bigr)$:
$203.771 - 63.080x = 167.040 - 12.384x$, logo $x = 36.731/50.696 = 0.7245$: a
\omterm{def:b1:rock-forming-minerals:groups}{olivina} é $\text{Fo}_{72}$ (arredondado ao percentual mais próximo).
\end{solution}

\begin{solution}{exo:b1:rock-forming-minerals:10}
Com $M_i = \rho_i V_i$,
\[
  \rho(x) - \bigl(x\rho_1 + (1 - x)\rho_2\bigr)
  = \frac{x\rho_1V_1 + (1 - x)\rho_2V_2 - \bigl(x\rho_1 + (1 - x)\rho_2\bigr)\bigl(xV_1 + (1 - x)V_2\bigr)}{V(x)} .
\]
Desenvolvendo o produto, o numerador é $x\rho_1V_1 - x^2\rho_1V_1 +
(1 - x)\rho_2V_2 - (1 - x)^2\rho_2V_2 - x(1 - x)(\rho_1V_2 + \rho_2V_1)$, e
$x - x^2 = x(1 - x)$, $(1 - x) - (1 - x)^2 = x(1 - x)$, logo ele é
$x(1 - x)(\rho_1V_1 + \rho_2V_2 - \rho_1V_2 - \rho_2V_1) =
x(1 - x)(\rho_1 - \rho_2)(V_1 - V_2)$. Para a \omterm{def:b1:rock-forming-minerals:groups}{olivina} (1 = forsterita) as duas
diferenças são negativas, $3.275 - 4.392 = -1.117$ e $42.96 - 46.40 =
-3.44$, de modo que o produto é positivo e a densidade fica acima da reta. Em $\text{Fo}_{50}$: $0.25\times1.117\times3.44/44.68 =
\qty{0.0215}{g.cm^{-3}}$, o que concorda com \cref{exo:b1:rock-forming-minerals:9}, dentro do arredondamento dos volumes molares.
\end{solution}

\begin{solution}{exo:b1:rock-forming-minerals:11}
Total \qty{99.60}{\percent}. Moles de óxido: \ce{SiO2} $0.90042$,
\ce{Al2O3} $0.28344$, \ce{CaO} $0.19973$, \ce{Na2O} $0.08229$,
\ce{K2O} $0.00318$. Oxigênios: $1.80084 + 0.85033 + 0.19973 + 0.08229 +
0.00318 = 2.93636$; $f = 8/2.93636 = 2.72446$. Cátions (o dobro dos mols para
\ce{Al2O3}, \ce{Na2O}, \ce{K2O}): Si \num{2.453}, Al \num{1.544}, Ca
\num{0.544}, Na \num{0.448}, K \num{0.017}. Sítios tetraédricos:
$2.453 + 1.544 = 3.997$; sítio grande: $0.544 + 0.448 + 0.017 = 1.009$. Então
$\mathrm{Al} - 1 = 0.544$, igual ao Ca dentro do arredondamento: cada cálcio entrou com um alumínio. $\mathrm{An} = 100\times0.544/1.009 = 53.9$: uma labradorita, $\text{An}_{54}$.
\end{solution}

\begin{solution}{exo:b1:rock-forming-minerals:12}
(a) $M(\ce{Fe3O4}) = 3\times55.845 + 4\times15.999 = \qty{231.531}{g.mol^{-1}}$.
Seus três ferros, registrados como \ce{FeO}, pesam $3\times71.844 = \qty{215.532}{g}$,
isto é, $215.532/231.531 = \qty{93.09}{\percent}$; o oxigênio ligado ao
\ce{Fe^{3+}}, além do que \ce{FeO} contabiliza, falta no total.
(b) $93.09/71.844 = 1.29572$ mol de \ce{FeO} e de oxigênio; $f =
4/1.29572 = 3.08708$; quatro ferros por quatro oxigênios, \ce{Fe4O4}. A magnetita tem
três cátions para quatro oxigênios. Os oxigênios foram subcontados: cada
\ce{Fe^{3+}} tem um oxigênio e meio, e não um, como supõe \ce{FeO}. (c) Em três cátions, Fe $= 3$
com $a$ \ce{Fe^{2+}} e $b$ \ce{Fe^{3+}}: $a + b = 3$ e $2a + 3b = 8$,
logo $b = 2$, $a = 1$: $\mathrm{Fe^{2+}Fe^{3+}_2O_4}$. (d) Um \ce{FeO} e
um \ce{Fe2O3} por unidade de fórmula: $71.844/231.531 = \qty{31.03}{\percent}$
de \ce{FeO} e $159.687/231.531 = \qty{68.97}{\percent}$ de \ce{Fe2O3}, total
\qty{100.00}{\percent}.
\end{solution}

\begin{solution}{pb:b1:rock-forming-minerals:1}
\textbf{1.} \ce{SiO2} $28.085 + 2\times15.999 = 60.083$; \ce{FeO} 71.844;
\ce{MnO} 70.937; \ce{MgO} 40.304; \ce{NiO} 74.692; \ce{CaO} 56.077
(\si{g.mol^{-1}}).

\textbf{2.} $39.50 + 16.53 + 0.23 + 43.32 + 0.20 + 0.22 =
\qty{100.00}{\percent}$. Um total de \qty{97}{\percent} teria sugerido um componente não analisado (água, ou oxigênio ligado a \ce{Fe^{3+}} registrado como \ce{FeO}), ou uma análise ruim.

\textbf{3.} \ce{SiO2} $39.50/60.083 = 0.65742$; \ce{FeO} 0.23008; \ce{MnO}
0.00324; \ce{MgO} 1.07483; \ce{NiO} 0.00268; \ce{CaO} 0.00392.

\textbf{4.} $2\times0.65742 = 1.31485$ para \ce{SiO2}, os mesmos números da
questão 3 para os demais; soma $1.31485 + 0.23008 + 0.00324 + 1.07483 +
0.00268 + 0.00392 = 2.62960$.

\textbf{5.} Cada óxido tem um cátion: os números da questão 3.

\textbf{6.} $f = 4/2.62960 = 1.52114$.

\textbf{7.} Si $0.65742\times1.52114 = 1.000$; Fe 0.350; Mn 0.005; Mg 1.635;
Ni 0.004; Ca 0.006.

\textbf{8.} $1.635 + 0.350 + 0.005 + 0.004 + 0.006 = 2.000$ para os
sítios octaédricos e 1.000 para o sítio tetraédrico: exatamente os 2 e 1 ideais.

\textbf{9.} \ce{(Mg_{1.635}Fe_{0.350}Mn_{0.005}Ni_{0.004}Ca_{0.006})SiO4}.

\textbf{10.} $2\times2.000 + 4\times1.000 = +8.000$ contra $4\times(-2) =
-8$. Cada cátion foi contado com a carga do seu óxido, de modo que as cargas positivas são o dobro dos oxigênios por construção: a verificação não pode falhar. Os totais de sítios da questão 8 poderiam ter falhado.

\textbf{11.} O níquel difere do magnésio por
$(0.69 - 0.72)/0.72 = -\qty{4}{\percent}$, bem dentro de \qty{15}{\percent}:
ele substitui livremente. O cálcio é $(1.00 - 0.72)/0.72 = \qty{39}{\percent}$
maior, além de \qty{30}{\percent}: quase não entra.

\textbf{12.} Ferro contado como \ce{FeO} quando parte dele é \ce{Fe^{3+}} deixa
poucos oxigênios por cátion, daí mais de três cátions em quatro oxigênios
(\cref{exo:b1:rock-forming-minerals:12}). A soma é $1.000 + 2.000 = 3.000$: sem \ce{Fe^{3+}} significativo.

\textbf{13.} $\mathrm{Fo} = 100\times1.635/(1.635 + 0.350) = 163.5/1.985 =
82.4$.

\textbf{14.} $\mathrm{Mg\#} = 82.4$: todo o ferro é \ce{Fe^{2+}}
(questão 12), de modo que $\mathrm{Fe^{2+}} = \mathrm{Fe}$ e as duas definições
coincidem.

\textbf{15.} $\text{Fo}_{82.4}\text{Fa}_{17.6}$, $x = 0.824$; o ferro ocupa
\qty{17.6}{\percent} dos sítios octaédricos no modelo binário
($0.350/2.000 = \qty{17.5}{\percent}$ com os elementos menores).

\textbf{16.} O fator multiplica Mg e Fe do mesmo modo e se cancela:
$\mathrm{Fo} = 100\times1.07483/(1.07483 + 0.23008) = 82.4$.

\textbf{17.} $40.5/40.304 = 1.00486$ e $20.0/71.844 = 0.27838$;
$\mathrm{Fo} = 100\times1.00486/1.28324 = 78.3$. A \omterm{def:b1:rock-forming-minerals:groups}{olivina} posterior é mais rica
em ferro.

\textbf{18.} $M = 140.691$ e \qty{203.771}{g.mol^{-1}}; os volumes molares são
$V = 140.691/3.275 = 42.96$ e $203.771/4.392 =
\qty{46.40}{cm^3.mol^{-1}}$.

\textbf{19.} $M = 0.824\times140.691 + 0.176\times203.771 = 115.929 +
35.864 = \qty{151.793}{g.mol^{-1}}$; $V = 0.824\times42.96 +
0.176\times46.40 = 35.399 + 8.166 = \qty{43.565}{cm^3.mol^{-1}}$.

\textbf{20.} $\rho = 151.793/43.565 = \qty{3.484}{g.cm^{-3}}$.

\textbf{21.} $0.824\times3.275 + 0.176\times4.392 = 2.699 + 0.773 =
\qty{3.472}{g.cm^{-3}}$, \qty{0.012}{g.cm^{-3}} (\qty{0.3}{\percent}) abaixo do valor da mistura ideal: próximo, como prevê \cref{exo:b1:rock-forming-minerals:10}.

\textbf{22.} $0.12\times\qty{3484}{kg.m^{-3}}\times\qty{1}{m^3} =
\qty{418}{kg}$.

\textbf{23.} $418/2950 = \qty{14.2}{\percent}$: a \omterm{def:b1:rock-forming-minerals:groups}{olivina} densa compõe uma
parcela maior da massa do que do volume.

\textbf{24.} A \omterm{def:b1:rock-forming-minerals:groups}{olivina} é $\text{Fo}_{82}$: seu teor de forsterita é
\textbf{82 por cento}.
\end{solution}
